Indeed, if @@M0@@, @@M1@@ is evidently an @@M2@@-sheaf. If @@M3@@, by 2.4, @@M4@@ since @@M5@@. Let @@M6@@ be a hyperplane section of @@M7@@ and let @@M8@@ be an exact sequence defined by @@M9@@. If @@M10@@, by 1.7, @@M11@@; thus @@M12@@ is a @@M13@@-sheaf. If @@M14@@, one can take @@M15@@ integral by Bertini's theorem and one may thus suppose by induction that @@M16@@ is a @@M17@@-sheaf; this gives the assertion.

\sgaStableTarget{sga6.e13.lem.2-6}\textbf{Lemma 2.7.} \emph{If @@M18@@ and @@M19@@, then}:
\begin{enumerate}[label=(\roman*)]
\item \emph{For every @@M20@@, @@M21@@}
\item \emph{@@M22@@.}
\end{enumerate}

Let @@M23@@ be an integral hyperplane section of @@M24@@. By 1.7 and 2.5, @@M25@@ is a @@M26@@-sheaf with Hilbert polynomial @@M27@@. Thus, by 1.11, @@M28@@ and @@M29@@ is @@M30@@-regular.

Consider the following exact sequence:
@@M31@@
For @@M32@@, we have @@M33@@ since @@M34@@ by 1.3 (i); for @@M35@@, we have @@M36@@ since @@M37@@ by 2.4. Finally, for every @@M38@@, @@M39@@ since there exists an injection @@M40@@, and @@M41@@ since @@M42@@; thus, for @@M43@@, we have @@M44@@. As a consequence, we conclude that for every @@M45@@,
\sgaStableTarget{sga6.e13.eq.2-6-1}@@M46@@

